\chapter{转牌 / Turn}

\section{本章导入}

本章讨论转牌圈。转牌的出现会显著改变局势：一些牌会增强进攻方范围，一些牌会补全防守方听牌，也有一些牌会让原本清晰的价值牌变得尴尬。转牌不是翻牌策略的机械延续，而是一次重新评估。

学习转牌的目标，是理解哪些转牌适合继续进攻，哪些转牌适合控池，哪些局面应当承认计划失败并放弃。

\section{核心概念}

本章将介绍好转牌、坏转牌、范围变化、第二枪、控池、延迟价值下注和放弃策略。读者需要看到，转牌决策往往比翻牌更考验计划性。

\section{新手常见误区}

新手常在转牌犯两类错误：一类是翻牌下注后无论转牌如何都继续下注，另一类是一遇到压力就放弃所有中等牌力。本节后续将讨论如何在二者之间找到依据。

\section{实战思考框架}

转牌应按“转牌是否改变坚果结构--谁的范围更受益--我的手牌功能是否变化--下注后河牌计划是什么”的顺序思考。没有河牌计划的转牌下注，往往会制造更昂贵的错误。

\section{牌例拆解}

\subsection{牌例一：待补充}

本牌例将用于分析进攻方在有利转牌上的第二枪。

\subsection{牌例二：待补充}

本牌例将用于分析防守方在危险转牌上的跟注与弃牌边界。

\section{本章总结}

转牌是局势变化最明显的一街。读者需要学会根据出张重新评估范围和牌力，而不是依赖翻牌圈的第一印象。

\section{课后训练}

请找出三个你曾经在转牌犹豫的牌局，分别判断转牌属于好牌、坏牌还是中性牌，并说明原因。
